\documentclass[a4paper,11pt]{article}
\usepackage[a4paper, margin=8em]{geometry}

% usa i pacchetti per la scrittura in italiano
\usepackage[french,italian]{babel}
\usepackage[T1]{fontenc}
\usepackage[utf8]{inputenc}
\frenchspacing 

% usa i pacchetti per la formattazione matematica
\usepackage{amsmath, amssymb, amsthm, amsfonts}

% usa altri pacchetti
\usepackage{gensymb}
\usepackage{hyperref}
\usepackage{standalone}

\usepackage{colortbl}

\usepackage{xstring}
\usepackage{karnaugh-map}

% imposta il titolo
\title{Notes on Liu-Layland bounds for RM scheduling}
\author{Luca Seggiani}
\date{2026}

% spaziatura liste
\usepackage{enumitem}
\setlist[itemize]{noitemsep, topsep=4pt, partopsep=0pt, parsep=2pt, leftmargin=*}
\setlist[enumerate]{noitemsep, topsep=4pt, partopsep=0pt, parsep=2pt, leftmargin=*}

% imposta lo stile
% usa helvetica
\usepackage[scaled]{helvet}
% usa palatino
\usepackage{palatino}
% usa un font monospazio guardabile
\usepackage{lmodern}

\renewcommand{\rmdefault}{ppl}
\renewcommand{\sfdefault}{phv}
\renewcommand{\ttdefault}{lmtt}

% circuiti
\usepackage{circuitikz}
\usetikzlibrary{babel}

% testo cerchiato
\newcommand*\circled[1]{\tikz[baseline=(char.base)]{
            \node[shape=circle,draw,inner sep=2pt] (char) {#1};}}

% disponi il titolo
\makeatletter
\renewcommand{\maketitle} {
	\begin{center} 
		\begin{minipage}[t]{.8\textwidth}
			\textsf{\huge\bfseries \@title} 
		\end{minipage}%
		\begin{minipage}[t]{.2\textwidth}
			\raggedleft \vspace{-1.65em}
			\textsf{\small \@author} \vfill
			\textsf{\small \@date}
		\end{minipage}
		\par
	\end{center}

	\thispagestyle{empty}
	\pagestyle{fancy}
}
\makeatother

% disponi teoremi
\usepackage{tcolorbox}
\newtcolorbox[auto counter, number within=section]{theorem}[2][]{%
	colback=blue!10, 
	colframe=blue!40!black, 
	sharp corners=northwest,
	fonttitle=\sffamily\bfseries, 
	title=Theorem~\thetcbcounter: #2, 
	#1
}

% disponi definizioni
\newtcolorbox[auto counter, number within=section]{definition}[2][]{%
	colback=red!10,
	colframe=red!40!black,
	sharp corners=northwest,
	fonttitle=\sffamily\bfseries,
	title=Definition~\thetcbcounter: #2,
	#1
}

% disponi codice
\usepackage{listings}
\usepackage[table]{xcolor}

\definecolor{codegreen}{rgb}{0,0.6,0}
\definecolor{codegray}{rgb}{0.5,0.5,0.5}
\definecolor{codepurple}{rgb}{0.58,0,0.82}
\definecolor{backcolour}{rgb}{0.95,0.95,0.92}

\lstdefinestyle{codestyle}{
		backgroundcolor=\color{black!5}, 
		commentstyle=\color{codegreen},
		keywordstyle=\bfseries\color{magenta},
		numberstyle=\sffamily\tiny\color{black!60},
		stringstyle=\color{green!50!black},
		basicstyle=\ttfamily\footnotesize,
		breakatwhitespace=false,         
		breaklines=true,                 
		captionpos=b,                    
		keepspaces=true,                 
		numbers=left,                    
		numbersep=5pt,                  
		showspaces=false,                
		showstringspaces=false,
		showtabs=false,                  
		tabsize=2
}

\lstdefinestyle{shellstyle}{
		backgroundcolor=\color{black!5}, 
		basicstyle=\ttfamily\footnotesize\color{black}, 
		commentstyle=\color{black}, 
		keywordstyle=\color{black},
		numberstyle=\color{black!5},
		stringstyle=\color{black}, 
		showspaces=false,
		showstringspaces=false, 
		showtabs=false, 
		tabsize=2, 
		numbers=none, 
		breaklines=true
}


\lstdefinelanguage{x86}{ 
  keywords={AAA, AAD, AAM, AAS, ADC, ADCB, ADCW, ADCL, ADD, ADDB, ADDW, ADDL, AND, ANDB, ANDW, ANDL,
        ARPL, BOUND, BSF, BSFL, BSFW, BSR, BSRL, BSRW, BSWAP, BT, BTC, BTCB, BTCW, BTCL, BTR, 
        BTRB, BTRW, BTRL, BTS, BTSB, BTSW, BTSL, CALL, CBW, CDQ, CLC, CLD, CLI, CLTS, CMC, CMP,
        CMPB, CMPW, CMPL, CMPS, CMPSB, CMPSD, CMPSW, CMPXCHG, CMPXCHGB, CMPXCHGW, CMPXCHGL,
        CMPXCHG8B, CPUID, CWDE, DAA, DAS, DEC, DECB, DECW, DECL, DIV, DIVB, DIVW, DIVL, ENTER,
        HLT, IDIV, IDIVB, IDIVW, IDIVL, IMUL, IMULB, IMULW, IMULL, IN, INB, INW, INL, INC, INCB,
        INCW, INCL, INS, INSB, INSD, INSW, INT, INT3, INTO, INVD, INVLPG, IRET, IRETD, JA, JAE,
        JB, JBE, JC, JCXZ, JE, JECXZ, JG, JGE, JL, JLE, JMP, JNA, JNAE, JNB, JNBE, JNC, JNE, JNG,
        JNGE, JNL, JNLE, JNO, JNP, JNS, JNZ, JO, JP, JPE, JPO, JS, JZ, LAHF, LAR, LCALL, LDS,
        LEA, LEAVE, LES, LFS, LGDT, LGS, LIDT, LMSW, LOCK, LODSB, LODSD, LODSW, LOOP, LOOPE,
        LOOPNE, LSL, LSS, LTR, MOV, MOVB, MOVW, MOVL, MOVSB, MOVSD, MOVSW, MOVSX, MOVSXB,
        MOVSXW, MOVSXL, MOVZX, MOVZXB, MOVZXW, MOVZXL, MUL, MULB, MULW, MULL, NEG, NEGB, NEGW,
        NEGL, NOP, NOT, NOTB, NOTW, NOTL, OR, ORB, ORW, ORL, OUT, OUTB, OUTW, OUTL, OUTSB, OUTSD,
        OUTSW, POP, POPL, POPW, POPB, POPA, POPAD, POPF, POPFD, PUSH, PUSHL, PUSHW, PUSHB, PUSHA, 
        RCLL, RCR, RCRB, RCRW, RCRL, RDMSR, RDPMC, RDTSC, REP, REPE, REPNE, RET, ROL, ROLB, ROLW,
        ROLL, ROR, RORB, RORW, RORL, SAHF, SAL, SALB, SALW, SALL, SAR, SARB, SARW, SARL, SBB,
        SBBB, SBBW, SBBL, SCASB, SCASD, SCASW, SETA, SETAE, SETB, SETBE, SETC, SETE, SETG, SETGE,
        SETL, SETLE, SETNA, SETNAE, SETNB, SETNBE, SETNC, SETNE, SETNG, SETNGE, SETNL, SETNLE,
        SETNO, SETNP, SETNS, SETNZ, SETO, SETP, SETPE, SETPO, SETS, SETZ, SGDT, SHL, SHLB, SHLW,
        SHLL, SHLD, SHR, SHRB, SHRW, SHRL, SHRD, SIDT, SLDT, SMSW, STC, STD, STI, STOSB, STOSD,
        STOSW, STR, SUB, SUBB, SUBW, SUBL, TEST, TESTB, TESTW, TESTL, VERR, VERW, WAIT, WBINVD,
        XADD, XADDB, XADDW, XADDL, XCHG, XCHGB, XCHGW, XCHGL, XLAT, XLATB, XOR, XORB, XORW, XORL},
  keywordstyle=\color{blue}\bfseries,
  ndkeywordstyle=\color{darkgray}\bfseries,
  identifierstyle=\color{black},
  sensitive=false,
  comment=[l]{\#},
  morecomment=[s]{/*}{*/},
  commentstyle=\color{purple}\ttfamily,
  stringstyle=\color{red}\ttfamily,
  morestring=[b]',
  morestring=[b]"
}

\lstdefinelanguage{riscv}{
  keywords={
    LUI, AUIPC,
    JAL, JALR,
    BEQ, BNE, BLT, BGE, BLTU, BGEU,
    LB, LH, LW, LBU, LHU, LD,
    SB, SH, SW, SD,
    ADDI, SLTI, SLTIU, XORI, ORI, ANDI,
    SLLI, SRLI, SRAI,
    ADD, SUB, SLL, SLT, SLTU, XOR, SRL, SRA, OR, AND,
    ADDIW, SLLIW, SRLIW, SRAIW,
    ADDW, SUBW, SLLW, SRLW, SRAW,
    MUL, MULH, MULHSU, MULHU, DIV, DIVU, REM, REMU,
    MULW, DIVW, DIVUW, REMW, REMUW,
    LR.W, SC.W, AMOSWAP.W, AMOADD.W, AMOXOR.W,
    AMOAND.W, AMOOR.W, AMOMIN.W, AMOMAX.W,
    AMOMINU.W, AMOMAXU.W,
    LR.D, SC.D, AMOSWAP.D, AMOADD.D, AMOXOR.D,
    AMOAND.D, AMOOR.D, AMOMIN.D, AMOMAX.D,
    AMOMINU.D, AMOMAXU.D,
    FLW, FSW,
    FMADD.S, FMSUB.S, FNMSUB.S, FNMADD.S,
    FADD.S, FSUB.S, FMUL.S, FDIV.S, FSQRT.S,
    FSGNJ.S, FSGNJN.S, FSGNJX.S,
    FMIN.S, FMAX.S,
    FCVT.W.S, FCVT.WU.S, FCVT.S.W, FCVT.S.WU,
    FMV.X.W, FMV.W.X,
    FEQ.S, FLT.S, FLE.S,
    FLD, FSD,
    FMADD.D, FMSUB.D, FNMSUB.D, FNMADD.D,
    FADD.D, FSUB.D, FMUL.D, FDIV.D, FSQRT.D,
    FSGNJ.D, FSGNJN.D, FSGNJX.D,
    FMIN.D, FMAX.D,
    FCVT.W.D, FCVT.WU.D, FCVT.D.W, FCVT.D.WU,
    FCVT.S.D, FCVT.D.S,
    FMV.X.D, FMV.D.X,
    FEQ.D, FLT.D, FLE.D,
    CSRRW, CSRRS, CSRRC,
    CSRRWI, CSRRSI, CSRRCI,
    FENCE.I,
    FENCE,
    ECALL, EBREAK,
    MRET, SRET, URET,
    WFI,
    SFENCE.VMA,
    NOP, LI, LA, MV,
    NOT, NEG, NEGW,
    SEQZ, SNEZ, SLTZ, SGTZ,
    BEQZ, BNEZ, BLEZ, BGEZ, BLTZ, BGTZ,
    BGT, BLE, BGTU, BLEU,
    J, JR, RET,
    CALL, TAIL,
    RDINSTRET, RDCYCLE, RDTIME,
    CSRR, CSRW, CSRS, CSRC,
    CSRWI, CSRSI, CSRCI
  },
  keywordstyle=\color{blue}\bfseries,
  ndkeywordstyle=\color{darkgray}\bfseries,
  identifierstyle=\color{black},
  sensitive=false,
  comment=[l]{\#},
  morecomment=[s]{/*}{*/},
  commentstyle=\color{purple}\ttfamily,
  stringstyle=\color{red}\ttfamily,
  morestring=[b]',
  morestring=[b]"
}

\lstdefinelanguage{arm}{
  keywords={
    ADC, ADD, AND, BIC, CMN, CMP, EOR, MOV, MVN,
    RSB, RSC, SBC, SUB, TST, TEQ,
    MLA, MLS, MUL, SMLAL, SMULL, UMLAL, UMULL,
    QADD, QSUB, QDADD, QDSUB,
    QADD16, QADD8, QASX, QSAX,
    QSUB16, QSUB8, QSAX, QSUBADDX,
    ASR, LSL, LSR, ROR, RRX,
    B, BL, BX, BLX,
    BXJ,
    BEQ, BNE, BCS, BHS, BCC, BLO,
    BMI, BPL, BVS, BVC,
    BHI, BLS, BGE, BLT, BGT, BLE,
    BAL,
    LDR, LDRB, LDRH, LDRSB, LDRSH,
    LDRD, LDM, LDMIA, LDMIB, LDMDA, LDMDB,
    LDRT, LDRBT, LDRHT, LDRSBT, LDRSHT,
    STR, STRB, STRH, STRD,
    STM, STMIA, STMIB, STMDA, STMDB,
    STRT, STRBT, STRHT,
    PUSH, POP,
    SWP, SWPB,
    REV, REV16, REVSH,
    BFC, BFI, BFXIL, SBFX, UBFX,
    CLZ,
    LDREX, LDREXB, LDREXH, LDREXD,
    STREX, STREXB, STREXH, STREXD,
    DMB, DSB, ISB,
    MCR, MCRR, MRC, MRRC,
    MRS, MSR,
    CPS, SETEND, SVC, BKPT,
    CLREX, WFE, WFI, SEV,
    YIELD, NOP,
    VLDR, VSTR,
    VMOV, VABS, VNEG,
    VADD, VSUB, VMUL, VDIV,
    VMLA, VMLS, VNMLA, VNMLS, VNMUL,
    VMAX, VMIN,
    VSQRT,
    VCMP, VCMPE,
    VCVT,
    VCLZ, VCLS,
    VAND, VORR, VEOR, VBIC, VMVN,
    VLD1, VLD2, VLD3, VLD4,
    VST1, VST2, VST3, VST4,
    ADR, ADCS, ADDS, SUBS, MOVS, ANDS, ORRS, EORS
  },
  keywordstyle=\color{blue}\bfseries,
  ndkeywordstyle=\color{darkgray}\bfseries,
  identifierstyle=\color{black},
  sensitive=false,
  comment=[l]{@},
  morecomment=[l]{;},
  morecomment=[s]{/*}{*/},
  commentstyle=\color{purple}\ttfamily,
  stringstyle=\color{red}\ttfamily,
  morestring=[b]',
  morestring=[b]"
}

\lstset{style=codestyle, language=C}

% disponi sezioni
\usepackage{titlesec}

\titleformat{\section}
	{\sffamily\Large\bfseries} 
	{\thesection}{1em}{} 
\titleformat{\subsection}
	{\sffamily\large\bfseries}   
	{\thesubsection}{1em}{} 
\titleformat{\subsection}
	{\sffamily\normalsize\bfseries} 
	{\thesubsection}{1em}{}

% tikz
\usepackage{tikz}

% float
\usepackage{float}

% grafici
\usepackage{pgfplots}
\pgfplotsset{width=10cm,compat=1.9}

% disponi alberi
\usepackage{forest}

\forestset{
	rectstyle/.style={
		for tree={rectangle,draw,font=\large\sffamily}
	},
	roundstyle/.style={
		for tree={circle,draw,font=\large}
	}
}

% disponi algoritmi
\usepackage{algorithm}
\usepackage{algorithmic}
\makeatletter
\renewcommand{\ALG@name}{Algorithm}
\makeatother

% disponi numeri di pagina
\usepackage{fancyhdr}
\fancyhf{} 
\fancyfoot[L]{\sffamily{\thepage}}

\makeatletter
\fancyhead[L]{\raisebox{1ex}[0pt][0pt]{\sffamily{\@title \ \@date}}} 
\fancyhead[R]{\raisebox{1ex}[0pt][0pt]{\sffamily{\@author}}}
\makeatother

\begin{document}

\pagestyle{fancy}
\thispagestyle{empty}
\renewcommand{\thispagestyle}[1]{}

\maketitle

\section{Rate Monotonic scheduling bounds}

This note examines the derivation of the Liu–Layland utilization bound for \textbf{RM} (\textit{Rate Monotonic}) scheduling, starting with the two-task case and extending the argument to an arbitrary number of tasks. The goal is to understand how task periods and execution times constrain the maximum processor utilization that can be guaranteed schedulable under a fixed-priority assignment. The derivation follows the original 1973 paper, making explicit the critical-instant argument, the minimization of the utilization bound, and the assumptions required at each step. Finally, the restriction on period ratios is removed, establishing the general bound
\[
U_{\mathrm{bound}}=m\left(2^{1/m}-1\right).
\]

The original article, for reference, is \textit{Scheduling Algorithms for Multiprogramming in a Hard-Real-Time Environment} (1973), by C. L. Liu and James W. Layland.

\subsection{Two-process case}

Let's first consider two processes, $\tau_1$ and $\tau_2$, with $T_1 < T_2$ (and thus $\tau_1$ at higher priority than $\tau_2$).
We seek to evaluate the maximum CPU utilization $U$, that it is the maximum service time the two processes can request.

Let's fix $C_1$ and leave $C_2$ free.
We consider the critical zone of $\tau_2$, that is a period $T_2$ over which $\tau_2$ executes interleaved with $\tau_1$.
To see the maximum value of $C_2$ (from which we'll derive the maximum utilization $U$) we have to consider two distinct situations:
\begin{enumerate}
\item $\tau_1$ is able to finish servicing its last request within the critical zone of $\tau_2$, before the deadline $T_2$ of $\tau_2$ arrives.

\begin{center}
\begin{tikzpicture}[scale=0.6]
  % Coordinates
  \coordinate (B) at (0.00, 1.25);
  \coordinate (C) at (0.00, 0.50);
  \coordinate (D) at (3.00, 0.50);
  \coordinate (E) at (3.00, 0.00);
  \coordinate (F) at (3.00, 1.00);
  \coordinate (K) at (0.00, 3.25);
  \coordinate (L) at (0.00, 2.50);
  \coordinate (M) at (2.00, 2.50);
  \coordinate (N) at (2.00, 2.00);
  \coordinate (O) at (2.00, 3.00);
  \coordinate (P) at (4.00, 2.50);
  \coordinate (Q) at (4.00, 2.00);
  \coordinate (R) at (4.00, 3.00);
  \coordinate (D_1) at (3.50, 0.50);
  \coordinate (G) at (4.50, 2.50);
  \coordinate (H) at (0.00, 2.00);
  \coordinate (I) at (0.00, 3.00);
  \coordinate (A) at (0.00, 0.00);
  \coordinate (J) at (0.00, 1.00);
  \coordinate (Y) at (2.00, 3.25);
  \coordinate (T) at (2.40, 5.00);
  \coordinate (U) at (2.00, 4.50);
  \coordinate (V) at (2.75, 4.50);
  \coordinate (W) at (2.75, 4.00);
  \coordinate (X) at (2.75, 5.00);
  \coordinate (Z) at (3.00, 1.25);
  \coordinate (AA) at (3.00, 5.50);
  \coordinate (AB) at (3.00, -0.50);
  \coordinate (S) at (2.00, 4.00);
  \coordinate (AC) at (2.00, 5.00);

  % Drawing
  \draw[gray, dashed] (AA) -- (AB);
  \draw (C) -- (D);
  \draw (E) -- (F);
  \draw (L) -- (M);
  \draw (N) -- (O);
  \draw (M) -- (P);
  \draw (Q) -- (R);
  \node[right] at (D_1) {$\dots$};
  \node[right] at (G) {$\dots$};
  \node[above=-2mm] at (K) {$\tau_1$};
  \draw (H) -- (I);
  \node[above=-2mm] at (B) {$\tau_2$};
  \draw (A) -- (J);
  \draw (J) -- (A);
  \node[above=-2mm] at (Y) {$T_1$};
  \draw (U) -- (V);
  \draw (W) -- (X);
  \node[above=-2mm] at (Z) {$T_2$};
  \node[above=-2mm] at (T) {$C_1$};
  \draw (S) -- (AC);
\end{tikzpicture}
\end{center}

This is expressed by the inequality:
$$
C_1 \leq T_2 - T_1 \left\lfloor \frac{T_2}{T_1} \right\rfloor
$$

We can calculate the maximum value of $C_2$ by taking the whole $T_2$ period and assuming the $\tau_1$ process is fine taking its time $C_1$ every time a request for it appears (that is, $\left\lceil \frac{T_2}{T_1} \right\rceil$ for a period of $T_1$) in $T_2$:
$$
{C_2}_\text{max}' = T_2 - C_1 \left\lceil \frac{T_2}{T_1} \right\rceil
$$

From this we compute the maximum utilization $U_\text{max}'$:
$$
U_\text{max}' = 1 + C_1 \left( \frac{1}{T_1} - \frac{1}{T_2} \left\lceil \frac{T_2}{T_1} \right\rceil \right)
$$

\item $\tau_1$ is not able to finish servicing its last request within the critical zone of $\tau_2$, before the deadline $T_2$ of $\tau_2$ arrives. This means that its $C_1$ execution time, issued at the last request within $T_2$, "bleeds" after the $T_2$ deadline.

\begin{center}
\begin{tikzpicture}[scale=0.6]
  % Coordinates
  \coordinate (B) at (0.00, 1.25);
  \coordinate (C) at (0.00, 0.50);
  \coordinate (D) at (3.00, 0.50);
  \coordinate (E) at (3.00, 0.00);
  \coordinate (F) at (3.00, 1.00);
  \coordinate (K) at (0.00, 3.25);
  \coordinate (L) at (0.00, 2.50);
  \coordinate (M) at (2.00, 2.50);
  \coordinate (N) at (2.00, 2.00);
  \coordinate (O) at (2.00, 3.00);
  \coordinate (P) at (4.00, 2.50);
  \coordinate (Q) at (4.00, 2.00);
  \coordinate (R) at (4.00, 3.00);
  \coordinate (D_1) at (3.50, 0.50);
  \coordinate (G) at (4.50, 2.50);
  \coordinate (H) at (0.00, 2.00);
  \coordinate (I) at (0.00, 3.00);
  \coordinate (A) at (0.00, 0.00);
  \coordinate (J) at (0.00, 1.00);
  \coordinate (Y) at (2.00, 3.25);
  \coordinate (T) at (2.70, 5.00);
  \coordinate (U) at (2.00, 4.50);
  \coordinate (V) at (3.25, 4.50);
  \coordinate (W) at (3.25, 4.00);
  \coordinate (X) at (3.25, 5.00);
  \coordinate (Z) at (3.00, 1.25);
  \coordinate (AA) at (3.00, 5.50);
  \coordinate (AB) at (3.00, -0.50);
  \coordinate (S) at (2.00, 4.00);
  \coordinate (AC) at (2.00, 5.00);

  % Drawing
  \draw[gray, dashed] (AA) -- (AB);
  \draw (C) -- (D);
  \draw (E) -- (F);
  \draw (L) -- (M);
  \draw (N) -- (O);
  \draw (M) -- (P);
  \draw (Q) -- (R);
  \node[right] at (D_1) {$\dots$};
  \node[right] at (G) {$\dots$};
  \node[above] at (K) {$\tau_1$};
  \draw (H) -- (I);
  \node[above] at (B) {$\tau_2$};
  \draw (A) -- (J);
  \draw (J) -- (A);
  \node[above] at (Y) {$T_1$};
  \draw (U) -- (V);
  \draw (W) -- (X);
  \node[above] at (Z) {$T_2$};
  \node[above] at (T) {$C_1$};
  \draw (S) -- (AC);
\end{tikzpicture}
\end{center}

This is expressed by the opposite of the inequality above:
$$
C_1 \geq T_2 - T_1 \left\lfloor \frac{T_2}{T_1} \right\rfloor
$$

We can calculate the maximum value of $C_2$ differently than above.
The idea is that, for each request apart from the last one, $\tau_1$ will execute in $C_1$ within a $T_1$ period, leaving $T_1 - C_1$ time free for $\tau_2$.
The last period that falls in $T_2$, as we said, is not going to be enough to allow $\tau_1$ to finish, and as its priority is higher than $\tau_2$ it's going to "steal away" all of the end of $\tau_2$'s period.
Mathematically:
$$
{C_2}_\text{max}'' = \left\lfloor \frac{T_2}{T_1} \right\rfloor \left( T_1 - C_1 \right)
$$

Again, from this we compute the maximum utilization $U_\text{max}''$:
$$
U_\text{max}'' = \frac{T_1}{T_2}\left\lfloor \frac{T_2}{T_1} \right\rfloor + C_1 \left( \frac{1}{T_1} - \frac{1}{T_2} \left\lfloor \frac{T_2}{T_1} \right\rfloor \right)
$$

\end{enumerate}

The two equations we got for $U_\text{max}'$ and $U_\text{max}''$ are respectively decreasing and increasing in $C_1$, that is the least maximum bound is given at the intersection:
$$
C_1 = T_2 - T_1 \left\lfloor \frac{T_2}{T_1} \right\rfloor
$$

Now the algebra gets fucked up fast so we first say:
$$
q = \frac{T_2}{T_1}
$$

And subsequently:
$$
\left\lfloor q \right\rfloor = f, \quad
\left\lceil q \right\rceil = c
$$

Now we compute the two bounds separately:
\begin{enumerate}
\item For $U_\text{max}'$ we have:
$$
\begin{aligned}
U_\text{max}' = 1 + \left( T_2 - T_1 \left\lfloor \frac{T_2}{T_1} \right\rfloor \right) \left( \frac{1}{T_1} - \frac{1}{T_2} \left\lceil \frac{T_2}{T_1} \right\rceil \right) &= 1 + q - \left\lceil q \right\rceil - \left\lfloor q \right\rfloor + \frac{1}{q} \left\lfloor q \right\rfloor \left\lceil q \right\rceil \\
&= 1 + q - c - f + \frac{1}{q} cf 
\end{aligned}
$$

\item For $U_\text{max}'$ we have:
$$
\begin{aligned}
U_\text{max}'' = \frac{T_1}{T_2}\left\lfloor \frac{T_2}{T_1} \right\rfloor + \left( T_2 - T_1 \left\lfloor \frac{T_2}{T_1} \right\rfloor \right) \left( \frac{1}{T_1} - \frac{1}{T_2} \left\lfloor \frac{T_2}{T_1} \right\rfloor \right) &= \frac{\left\lfloor q \right\rfloor}{q} \left( 1 + \left\lfloor q \right\rfloor \right) + q - 2 \left\lfloor q \right\rfloor \\
&= \frac{f}{q}(1 + f) + q - 2f
\end{aligned}
$$
\end{enumerate}

Ain't that something.
At this point we first need to prove that $U_\text{max}'$ and $U_\text{max}''$ are in fact equal.
Dealing with ceiling and floor functions is tricky so we split the case where $q \in \mathbb{Z}^+$ and $q \in \mathbb{Q}^+$ (we don't really care about negative numbers as these are time periods, and we need not evaluate all of $\mathbb{R}$ as these are ratios between integers, though I imagine in $\mathbb{R}$ the calculations are similar).

\begin{itemize}
\item In $\mathbb{Z}^+$ we can say $q = f = c$ and the equations simplify quickly to $1$:
$$
\begin{aligned}
U_\text{max}' &= 1 + q - q + \frac{1}{q} q^2 = 1 \\
U_\text{max}' &= \frac{q}{q}(1 + q) + q - 2q = 1
\end{aligned}
$$

\item In $\mathbb{Q}^+$ we can similarly say that $c = f + 1$, which leads to:
$$
\begin{aligned}
	U_\text{max}' = 1 + q - f - 1 - f + \frac{1}{q} \left( f (f + 1) \right) &= q - 2f + \frac{f}{q} \left( 1 + f \right) \\
	U_\text{max}'  &=   q - 2f + \frac{f}{q} ( 1 + f )
\end{aligned}
$$
\end{itemize}

Now that we have a function $U_\text{max}$ we have to minimize it to find the least maximum bound.
A few considerations:
\begin{itemize}
\item $q \geq 1$, because:
$$
q = \frac{T_2}{T_1} \geq 1 \implies T_2 > T_1
$$

Had it been otherwise $\tau_1$ and $\tau_2$ would swap priorities;
\item As we found out before, $U_\text{max} = 1$ for each $q \in \mathbb{Z}^+$;
\item For each point $q \in \{ \mathbb{Q}^+ \setminus \mathbb{Z}^+ \}$, $f$ is basically fixed, and $q$ is our free variable. We can ignore the discontinuities at $q \in \mathbb{Z}^+$ as we are going to find the lower bound is lower than $1$.
\end{itemize}

This being said, we derive considering $f$ as a parameter:
$$
\frac{d}{dq} \left( q - 2f + \frac{f}{q}(1 + f) \right) = 1 - \frac{f(1+f)}{q^2}
$$

The minimum is found by setting this to zero:
$$
1 - \frac{f(1+f)}{q^2} = 0 \implies q^* = \sqrt{f(1 + f)}
$$

Let's think about what this means.
Our function has a stationary point $q_i$ for each point which respects:
$$
q_i = \sqrt{f_i(1 + f_i)}
$$

With $f_i = 1, 2, ...$ as $q_1$ with $f_i = 1$ comes out to $\sqrt{2}$, just above our $q > 1$ bound.
By quick evaluation at the first few points we see that they are indeed below $1$, so all we really have to do is find out which one of these is the lowest.

By substituting back in the original expression we get:
$$
U_\text{max}^* = \sqrt{f(1 + f)} - 2f + \frac{f}{\sqrt{f(1 + f)}} (1 + f) = 2 \left( \sqrt{f(1 + f)} - f\right)
$$

Let's derive, this time with $f$ as the variable:
$$
\frac{d}{df} \left( 2 \left( \sqrt{f(1 + f)} - f\right) \right) = 2 \left( \frac{1 + 2f}{2 \sqrt{f(1 + f)}} - 1 \right)
$$

Now, if only this derivative was greater than $0$, we could say that the sequence $q_i$ given by $f_i = 1, 2, 3...$ is non decreasing and thus the minimum is found at $q_1$ with $f_i = 1$.
Luck be on our side, we get:
$$
\begin{aligned}
2 \left( \frac{1 + 2f}{2 \sqrt{f(1 + f)}} - 1 \right) &> 0 \\
\frac{1 + 2f}{2 \sqrt{f(1 + f)}} &> \frac{1}{2} \\
1 + 2f &> 1 \sqrt{f (1 + f)} \\
1 + 3 f ^2 + 3 f &> 0
\end{aligned}
$$

Which is definitely true for $f > 0$.
We are happy.

In conclusion, what we got is that our CPU utilization function $U_\text{max}$ has a least maximum bound in:
$$
q^* = \frac{T_2}{T_1} = \sqrt{2}
$$

Which we got by imposing $f_i = 1$ in the $q_i = \sqrt{f_i(1 + f_i)}$ expression.
At this point all that's left is evauating the least maximum utilization:
$$
U_\text{max}^* = \sqrt{2} - 2 + \frac{1}{\sqrt{2}} \cdot 2 = 2 \sqrt{2} - 2 = 2(2^{\frac{1}{2}} - 1)
$$

\subsection{Multi-process case}

This should be nothing more than a generalization of the above case.
The original article is followable so I'll stay concise.

First off we take a set of $m$ processes, ordered as usual with periods:
$$
T_1 < T_2 < ... < T_m
$$

So that the priorities are ordered in the opposite direction, with $T_1$ having the highest priority and $T_m$ the lowest priority.
\textbf{Note} that for now we make an assumption: the ratio between any two request periods must be less than $2$.

We wish to show that if $V_1, ..., V_m$ make a minimal bound of the CPU utilization $U$, it must follow that:
$$
C_i = T_{i + 1} - T_i, \quad i \neq m
$$

We show this by proving the case for $C_i$, and letting the rest stem from recursion.
The equality to prove thus becomes:
$$
C_1 = T_2 - T_1
$$

Two situations to consider:
\begin{itemize}
\item If we perturb $C_1$ negatively, that is say that there exists some $\Delta > 0$ such that:
$$
C_1 = T_2 - T_1 + \Delta
$$

Then we can easily build a new set of runtimes $C_i$ which fully utilize the processor:
$$
\begin{cases}
C_1' &= T_2 - T_1 \\
C_2' &= C_2 + \Delta \\
C_3' &= C_3 \\
	 &\vdots \\
C_m' &= C_m
\end{cases}
$$

Let's consider the differences in utilization between the two sets.
Only the first two elements change, so:
$$
U - U' = \frac{\Delta}{T_1} - \frac{\Delta}{T_2}
$$

Since $T_1 < T_2$, $\frac{1}{T_1} > \frac{1}{T_2}$, and $U > U'$.
This is the first contradiction.

\item If we perturb $C_1$ positively, that is say that there exists some $\Delta > 0$ such that:
$$
C_1 = T_2 - T_1 - \Delta
$$

Then we can easily build a new set of runtimes $C_i$ which fully utilize the processor:
$$
\begin{cases}
C_1'' &= T_2 - T_1 \\
C_2'' &= C_2 - 2 \Delta \\
C_3'' &= C_3 \\
	 &\vdots \\
C_m'' &= C_m
\end{cases}
$$

Let's consider the differences in utilization between the two sets.
Only the first two elements change, so:
$$
U - U'' = -\frac{\Delta}{T_1} + 2 \frac{\Delta}{T_2}
$$

Which also leads to $U > U''$.
This is the second contradiction.
\end{itemize}

It follows that indeed:
$$
C_1 = T_2 - T_1
$$

And by recursion:
$$
C_i = T_{i + 1} - T_i, \quad i \neq m
$$

We can express this more compactly in a number of ways, of varying usefulness.
\begin{itemize}
\item First off, we are interested in $C_m$. By the assumption made above (the ratio between any two request periods must be less than $2$) we can say that:
$$
C_m = T_m - 2 \sum_i^{m - 1} C_i
$$

At this point note that the sum $\sum_i^{m - 1} C_i$ telescopes, so that this reduces to:
$$
C_i = T_{i + 1} - T_i, \quad
C_m = 2 T_1 - T_m
$$

\item "To simplify the notation" the paper now gives:
$$
g_i = \frac{T_m - T_i}{T_i}
$$

We can use this to give alternative definitions of $C_i$ and $C_m$:
$$
C_i = g_i T_i - g_{i + 1} T_{i + 1}, \quad
C_m = T_m - 2 g_1 T_1
$$

\end{itemize}

From the last set of equations we can derive a workable expression for the utilization $U$:
$$
\begin{aligned}
U
&= \sum_{i=1}^{m-1}
\left[
g_i - g_{i+1}\left(\frac{T_{i+1}}{T_i}\right)
\right]
+ 1 - 2g_1\left(\frac{T_1}{T_m}\right)
\\[6pt]
&= \sum_{i=1}^{m-1}
\left[
g_i - g_{i+1}\frac{g_i+1}{g_{i+1}+1}
\right]
+ 1 - 2\frac{g_1}{g_m+1}
\\[6pt]
&= 1
+ g_m\left(\frac{g_1-1}{g_m+1}\right)
+ \sum_{i=2}^{m-1}
g_i\left(\frac{g_i-g_{i-1}}{g_i+1}\right)
\end{aligned}
$$

As before, we can find a minimum by deriving with respect to $g_j$:
$$
\begin{aligned}
\frac{d U}{d g_j}
&=
\frac{g_j^2+2g_j-g_{j-1}}{(g_j+1)^2}
-\frac{g_{j+1}}{g_{j+1}+1}
=0
\end{aligned}
$$

Which can be proven to give:
$$
g_j = 2^{\frac{m - j}{m}} - 1, \quad j = 0, 1, 2, ..., m-1
$$

At this point we only need to put the $g_j$s back in the original equation, getting:
$$
U = m (2^{\frac{1}{m}} - 1)
$$

Most of the algebra wasn't notated here because I imagine it is unfathomably boring.

\par\smallskip

We still have to lift that last restriction, that the ratio between any two request periods must be less than $2$. 
This was useful to even get an equation for $C_m$ in the first place.

The is that any task set with a period ratio of at least $2$ can be transformed into another fully utilizing task set with no greater utilization.

\end{document}
